\documentclass[aspectratio=169]{beamer}
\usetheme{Madrid}
\usecolortheme{whale}

\usepackage[utf8]{inputenc}
\usepackage[spanish]{babel}
\usepackage{amsmath,amssymb}
\usepackage{graphicx}
\usepackage{listings}
\usepackage{booktabs}

\lstset{
    language=Python,
    basicstyle=\ttfamily\footnotesize,
    keywordstyle=\color{blue}\bfseries,
    stringstyle=\color{red},
    commentstyle=\color{gray}\itshape,
    breaklines=true,
    showstringspaces=false
}

\title[INT210 - Sesión 0.8]{Sesión 0.8: Ingesta de Archivos Delimitados "From Scratch"}
\subtitle{Módulo 0.2: Estructuras de Datos Nativas en Python}
\author[Diplomado en Ciencia de Datos]{Cátedra de Inteligencia Artificial y Ciencia de Datos}
\institute[INT210]{Machine Learning From Scratch}
\date{Semana 2}

\begin{document}

\frame{\titlepage}

\begin{frame}{Contenido de la Sesión}
    \tableofcontents
\end{frame}

\section{El Archivero con Llave y Context Managers}

\begin{frame}{La Analogía: El Archivero con Llave de Seguridad}
    \begin{columns}
        \column{0.55\textwidth}
        \begin{itemize}
            \item \textbf{\texttt{with open()}:} Abre la gaveta, procesa los datos y garantiza el cierre seguro al finalizar (patrón RAII).
            \item \textbf{\texttt{csv.reader}:} Lector en flujo (*Streaming*) que decodifica campos delimitados sin cargar todo el archivo en RAM.
            \item \textbf{\texttt{namedtuple}:} Registros inmutables con acceso por nombre de atributo y consumo mínimo de memoria.
        \end{itemize}
        \column{0.45\textwidth}
        \begin{alertblock}{Prevención de Fugas de Recursos}
            Omitir el administrador de contexto puede causar \textit{File Descriptor Leaks} y bloqueos del sistema operativo bajo alta concurrencia.
        \end{alertblock}
    \end{columns}
\end{frame}

\section{Arquitectura de Ingesta Nativa en Python}

\begin{frame}[fragile]{Pipeline de Ingesta Segura (Joel Grus, Cap. 9)}
\begin{lstlisting}
from collections import namedtuple
import csv

Transaccion = namedtuple('Transaccion', ['id', 'monto', 'ip'])

def cargar_datos(ruta):
    registros = []
    with open(ruta, mode='r', encoding='utf-8', newline='') as f:
        lector = csv.DictReader(f)
        for fila in lector:
            # Casteo estricto de tipos de datos
            reg = Transaccion(fila['id'], float(fila['monto']), fila['ip'])
            registros.append(reg)
    return registros
\end{lstlisting}
\end{frame}

\section{Comparación de Desempeño: Nativo vs Pandas}

\begin{frame}{Sobrecarga de Memoria y Latencia}
    \begin{center}
    \begin{tabular}{lll}
        \toprule
        \textbf{Criterio} & \textbf{Python Nativo (\texttt{csv} + \texttt{namedtuple})} & \textbf{Pandas DataFrame} \\
        \midrule
        \textbf{Memoria Base} & $\approx 10 - 15$ MB & $\approx 50 - 120$ MB \\
        \textbf{Cold Start} & $< 15$ ms & $1.5 - 3.0$ s \\
        \textbf{Dependencias} & Cero (Librería Estándar) & NumPy, pytz, C-extensions \\
        \textbf{Ideal Para} & Microservicios Serverless & Análisis Exploratorio Complejo \\
        \bottomrule
    \end{tabular}
    \end{center}
\end{frame}

\begin{frame}{Resumen General del Módulo 0}
    \begin{itemize}
        \item \textbf{Semana 1:} Dominio de bifurcaciones, \texttt{match-case}, bucles deterministas e indeterminados y matrices de confusión.
        \item \textbf{Semana 2:} Álgebra lineal vectorial con listas/tuplas, tablas hash con diccionarios/defaultdict, conjuntos de alta velocidad e ingesta pura "From Scratch".
        \item \textbf{¡Felicidades! Tienes las bases sólidas para adentrarte en los modelos de Machine Learning matemáticos del Capítulo 1.}
    \end{itemize}
\end{frame}

\end{document}
